% Lean module: MonadFormal.Approximation
\section{Approximation by statistical simulation}
\label{sec:simulation}

Exact factorization asks whether a required distinction is preserved.
Statistical simulation asks a different question: can the retained
experiment reproduce the distributions of another experiment well enough
for bounded decisions? We use normalized total variation and a common
Markov kernel. This is an instance of experiment comparison, not a new
scalar definition of MONAD defect~\cite{carterGaussianApproximation}.

Let $X,Y,Z$ be nonempty measurable spaces. Write $\mathcal P(X)$ for
probability measures on $X$ and
\[
 \mathcal T(X)=\{f:X\to[0,1]:f\text{ is measurable}\},\qquad
 d_{\mathrm{TV}}(\mu,\nu)
 =\sup_{f\in\mathcal T(X)}\left|\int f\,d\mu-\int f\,d\nu\right|.
\]
All these integrals exist: $f$ is bounded and the measures are finite.
The normalization uses $[0,1]$ tests, not $[-1,1]$ tests. A randomization
$K:X\rightsquigarrow Y$ is a Markov kernel: $K(x,\cdot)$ is a probability
measure and $x\mapsto K(x,B)$ is measurable for every measurable $B$.
Denote its action on a measure by $K\mu$, its composition with $L$ by
$LK$, and its input relabeling along a measurable map $q:Z\to X$ by
$K\circ q$, where $(K\circ q)(z,\cdot)=K(q(z),\cdot)$.
Its action on a test is
\[
 (K^*f)(x)=\int f(y)\,K(x,dy).
\]
\leandef{MonadFormal.Approximation.Test,MonadFormal.Approximation.tv,MonadFormal.Approximation.Randomization,MonadFormal.Approximation.Randomization.apply,MonadFormal.Approximation.Randomization.comp,MonadFormal.Approximation.Randomization.comap,MonadFormal.Approximation.Randomization.pull}

\begin{theorem}[Variation bounds and Markov contraction]
\label{thm:variation-contraction}
Normalized total variation is symmetric, vanishes on the diagonal,
is at most one, and satisfies the triangle inequality. Moreover,
\[
 d_{\mathrm{TV}}(K\mu,K\nu)\le d_{\mathrm{TV}}(\mu,\nu).
\]
\end{theorem}
\begin{proof}
Every test expectation lies in $[0,1]$, so its difference has absolute
value at most one. Symmetry and diagonal vanishing follow from the same
properties of absolute value. For each test, insert the intermediate
expectation and use the triangle inequality; take the supremum.

For contraction, measurability of kernel integration and the bounds on
$f$ show that $K^*f\in\mathcal T(X)$. Kernel integration gives
\[
 \int f\,d(K\mu)=\int K^*f\,d\mu.
\]
Thus every difference tested after $K$ is already bounded by the
supremum over tests before $K$. Taking the supremum proves contraction.
Kernel composition likewise gives $(LK)\mu=L(K\mu)$.
\leanproof{MonadFormal.Approximation.Test.integrable,MonadFormal.Approximation.Test.integral_bounds,MonadFormal.Approximation.test_le_tv,MonadFormal.Approximation.tv_le_one,MonadFormal.Approximation.tv_self,MonadFormal.Approximation.tv_symm,MonadFormal.Approximation.tv_triangle,MonadFormal.Approximation.Randomization.integral_apply,MonadFormal.Approximation.Randomization.apply_comp,MonadFormal.Approximation.tv_processing}
\end{proof}

% Lean module: MonadFormal.Approximation.Events
\subsection{Event meaning of the error}

\begin{theorem}[Event characterization and separation]
\label{thm:variation-events}
For probability measures on any measurable space,
\[
 d_{\mathrm{TV}}(\mu,\nu)
   =\sup_{B\text{ measurable}}|\mu(B)-\nu(B)|,
 \qquad d_{\mathrm{TV}}(\mu,\nu)=0\iff\mu=\nu.
\]
\end{theorem}
\begin{proof}
Indicator functions are admitted tests, which proves that the event
supremum is no larger than $d_{\mathrm{TV}}$.
For the reverse inequality, fix $f\in\mathcal T(X)$. The layer-cake
formula gives
\[
 \int f\,d\mu=\int_{(0,1]}\mu\{x:f(x)\ge t\}\,dt,
 \qquad
 \int f\,d\nu=\int_{(0,1]}\nu\{x:f(x)\ge t\}\,dt.
\]
The tail probabilities are measurable, bounded functions of $t$.
Subtract these identities. The absolute value of each integrand
difference is bounded by the event supremum; the integration interval
has length one. Hence the test difference obeys the same bound.
Taking the supremum over $f$ proves equality. Zero distance then forces
agreement on every measurable event, which is equality of measures;
the converse is diagonal vanishing.
\leanproof{MonadFormal.Approximation.Test.indicator,MonadFormal.Approximation.eventTV,MonadFormal.Approximation.eventTV_le_tv,MonadFormal.Approximation.Test.tail,MonadFormal.Approximation.Test.integral_eq_integral_tail,MonadFormal.Approximation.tv_eq_eventTV,MonadFormal.Approximation.tv_eq_zero_iff}
\end{proof}

% Lean module: MonadFormal.Approximation.Simulation
\subsection{Uniform decision transfer and composition}

An experiment is a family $E=(E_\theta)_{\theta\in\Theta}$ of
probability measures on a common sample space. Let $E,F,G$ have sample
spaces $X,Y,Z$, respectively, and the same parameter set $\Theta$.
Define the uniform simulation error and directed deficiency by
\[
 e(K;E,F)=\sup_{\theta\in\Theta}
   d_{\mathrm{TV}}(KE_\theta,F_\theta),\qquad
 \delta_{\mathrm{sim}}(E,F)=\inf_{K:X\rightsquigarrow Y} e(K;E,F).
\]
The infimum is over one common kernel, independent of $\theta$.
The direction $E\to F$ means that data from $E$ are used to simulate
$F$. This is Markov-kernel deficiency; no identification with more
general abstract transitions is assumed. The formal definitions use
extended nonnegative reals, with empty infimum $\infty$ and empty
supremum zero. Nonempty sample spaces ensure a constant randomization
exists, so the deficiencies used here are finite.
\leandef{MonadFormal.Approximation.Experiment,MonadFormal.Approximation.simulationError,MonadFormal.Approximation.deficiency}

\begin{lemma}[Normalization and exact simulation]
$0\le\delta_{\mathrm{sim}}(E,F)\le1$. If a common kernel satisfies
$KE_\theta=F_\theta$ for every $\theta$, then
$\delta_{\mathrm{sim}}(E,F)=0$. In particular,
$\delta_{\mathrm{sim}}(E,E)=0$.
\end{lemma}
\begin{proof}
Every kernel's error is at most one. The infimum is bounded above by
the error of a constant probability kernel and is nonnegative.
An exact simulator has zero error, so its deficiency is zero. Use the
identity kernel for the diagonal case.
\leanproof{MonadFormal.Approximation.simulationError_le_one,MonadFormal.Approximation.deficiency_le_simulationError,MonadFormal.Approximation.deficiency_le_one,MonadFormal.Approximation.deficiency_of_exact,MonadFormal.Approximation.deficiency_self}
\end{proof}

Let $D:Y\rightsquigarrow A$ be a randomized decision rule and let
$\ell_\theta:A\to[0,1]$ be measurable for each $\theta$. Its risk is
\[
 R_F(D,\ell;\theta)=\int_A\ell_\theta(a)\,(DF_\theta)(da).
\]
\leandef{MonadFormal.Approximation.risk}

\begin{theorem}[Bounded-risk transfer]\label{thm:risk-transfer}
Every simulator $K$ satisfies, uniformly over decisions, bounded losses
and parameters,
\[
 |R_E(DK,\ell;\theta)-R_F(D,\ell;\theta)|\le e(K;E,F).
\]
Consequently, for every $\varepsilon>0$ a common simulator can be chosen
so that, for all $D,\ell,\theta$,
\[
 |R_E(DK,\ell;\theta)-R_F(D,\ell;\theta)|
 <\delta_{\mathrm{sim}}(E,F)+\varepsilon.
\]
\end{theorem}
\begin{proof}
The loss is an admitted test on $A$. Its expectation difference is
bounded by the total variation between $DKE_\theta$ and $DF_\theta$.
Markov contraction bounds this by
$d_{\mathrm{TV}}(KE_\theta,F_\theta)\le e(K;E,F)$.
Choose $K$ with error below
$\delta_{\mathrm{sim}}(E,F)+\varepsilon$ by the defining infimum.
This choice uses the simulation error, before the loss, decision rule
or parameter is selected. No minimizing kernel is assumed to exist.
\leanproof{MonadFormal.Approximation.pointwise_le_simulationError,MonadFormal.Approximation.risk_transfer,MonadFormal.Approximation.risk_transfer_real,MonadFormal.Approximation.deficiency_risk_transfer,MonadFormal.Approximation.deficiency_risk_transfer_epsilon}
\end{proof}

\begin{theorem}[Composition of simulation errors]
\label{thm:simulation-composition}
For common kernels $K:X\rightsquigarrow Y$ and $L:Y\rightsquigarrow Z$,
\[
 e(LK;E,G)\le e(K;E,F)+e(L;F,G),
 \qquad
 \delta_{\mathrm{sim}}(E,G)
 \le\delta_{\mathrm{sim}}(E,F)+\delta_{\mathrm{sim}}(F,G).
\]
\end{theorem}
\begin{proof}
At each parameter, insert $LF_\theta$ as an intermediate measure:
\begin{align*}
 d_{\mathrm{TV}}(LKE_\theta,G_\theta)
 &\le d_{\mathrm{TV}}(LKE_\theta,LF_\theta)
       +d_{\mathrm{TV}}(LF_\theta,G_\theta)\\
 &\le d_{\mathrm{TV}}(KE_\theta,F_\theta)
       +d_{\mathrm{TV}}(LF_\theta,G_\theta).
\end{align*}
Take the supremum over $\theta$. The composite kernel is a candidate
for $\delta_{\mathrm{sim}}(E,G)$. Taking independent infima over $K$
and $L$ gives the deficiency bound; addition preserves these infima
in the extended nonnegative reals. Attainment is unnecessary.
\leanproof{MonadFormal.Approximation.simulationError_comp,MonadFormal.Approximation.deficiency_triangle}
\end{proof}

These are forward guarantees, not a proof of the converse randomization
theorem. No unbounded-loss, computational-cost, recovery, or causal
intervention guarantee follows without its own hypotheses. A restricted
class of implementable simulators needs explicit access and resource
rules, including closure and cost under composition. Small pairwise
error also does not supply an equivalence relation; see the exact
nontransitivity witness in the next section.

